%% reading-text-booklet.tex
%% ---------------------------------------------------------------------------
%% The companion text booklet. Two ways of referencing a text are supported:
%%
%%   \begin{ibtextbody}              numbers every fifth line in the gutter
%%   \begin{ibtextbody}[paragraphs]  no line numbers; use \ibparnum{n} instead
%%
%% Whichever you choose must match the rubrics in the question booklet:
%% \iblines{1}{15} for line-numbered texts, \ibpara{4} / \ibparas{3}{4} for
%% paragraph-numbered ones.
%%
%% All three texts here are original, written for this template. Any resemblance
%% to a real publication, person or work is coincidental.
%% ---------------------------------------------------------------------------
\documentclass[11pt,a4paper]{article}
\usepackage{ibenglishb}

\ibsetsession{QDHS\,--\,2026}
\ibsetmaxmarks{0}          % a text booklet carries no marks

\begin{document}

% ═══════════════════════════════════════ COVER ═════════════════════════════
\ibcover
  {English B -- Higher level -- Paper 2 -- Reading comprehension}
  {8 May 2026}
  {1 h}
  {\begin{ibinstructions}{Text booklet -- Instructions to students}
     \item Do not open this booklet until instructed to do so.
     \item This booklet accompanies paper 2 reading comprehension.
   \end{ibinstructions}}
  {5 pages}
  {\ibsessioncode}

% ══════════════════════ TEXT A — paragraph-numbered ════════════════════════
\ibtextlabel{A}
\ibtexttitle{The library that lends almost anything}

\begin{ibtextbody}[paragraphs]
  \ibparnum{1}
  On a Tuesday morning in the eastern suburbs, a retired electrician called
  Farida Osman is explaining to a stranger how to use a wallpaper steamer. The
  two of them are standing in what used to be the reference section of the
  Brookvale public library. The shelves are still there. What sits on them has
  changed completely.

  \ibparnum{2}
  Brookvale is one of a growing number of libraries that lend objects as well
  as books. Members pay nothing beyond their ordinary library membership, and
  in return they can borrow a carpet cleaner, a pasta machine, a telescope or a
  set of camping chairs for up to a fortnight. The collection now runs to just
  over two thousand items, and almost all of them were donated by people living
  within a few streets of the building.

  \ibparnum{3}
  The idea came from the branch manager, Tomas Iversen, who noticed how often
  borrowers asked whether the library had anything beyond books. He began with
  eleven power tools kept behind the desk. Within a month they were permanently
  out on loan, and he had a waiting list. What convinced the council to fund
  the expansion was not the popularity of the scheme but a calculation: the
  eleven drills had replaced, by his estimate, more than sixty purchases.

  \ibparnum{4}
  Not everything has worked. An early attempt to lend musical instruments was
  abandoned after a term, partly because repairs proved expensive and partly
  because borrowers kept them far longer than the two-week limit. Iversen is
  candid about it. The instruments, he says, were the only items nobody wanted
  to give back, which taught him something useful about which objects belong in
  a lending collection and which do not.

  \ibparnum{5}
  Farida joined as a volunteer in the first year and now runs a Saturday
  session where members can bring in something broken and be shown how to mend
  it. She is careful to describe what she does as teaching rather than
  repairing. If she simply fixed things, she says, the same item would come
  back in six months. Roughly two-thirds of what arrives at her table leaves
  working, which she considers a better result than it sounds, since the
  alternative for most of it was the tip.

  \ibparnum{6}
  Eleven similar collections have opened in the region since Brookvale began,
  and Iversen now spends part of each month advising the librarians running
  them. His advice is consistent and slightly deflating: start small, refuse
  donations you cannot store, and accept that the collection will be shaped by
  what your members happen to own rather than by what you would like to offer.
\end{ibtextbody}

\ibglossnote{the tip}{a place where household waste is taken}

% ═════════════════════════ TEXT B — line-numbered ══════════════════════════
\ibtextlabel{B}
\ibtexttitle{Five ways to make a new habit stick}

\begin{ibtextbody}
  Every January, gyms fill up and language apps report a surge of new
  downloads. By March, most of those accounts are dormant again. The problem is
  rarely a lack of willingness. People who abandon a new habit are usually
  people who wanted it badly and went about it in a way that could not last.
  Here are five approaches that research and experience both support.

  \textbf{Make the first version absurdly small}\\
  The most common mistake is starting at the size you eventually want. Someone
  who has not run for a decade decides to run five kilometres, manages it
  twice, and stops. A better opening is one that feels almost too easy: two
  minutes of practice, a single page, one length of the pool. The purpose of
  the first fortnight is not progress. It is to establish that this is
  something you do.

  \ibgap{15}\\
  A habit attached to nothing in particular has to be remembered, and
  remembering is effortful. A habit attached to something you already do every
  day is prompted for you. Practise after you brush your teeth, or read while
  the kettle boils. The existing routine becomes the reminder, which means you
  are no longer relying on motivation at the precise moment motivation is
  lowest.

  \ibgap{16}\\
  Nearly everyone breaks a streak eventually, and what happens next matters far
  more than the break itself. People who treat a missed day as evidence of
  failure tend to abandon the habit within the week. People who treat it as an
  ordinary interruption simply resume. The useful rule is to never miss twice:
  one gap is an accident, two is the beginning of a new pattern.

  \ibgap{17}\\
  Habits maintained alone are fragile. Those maintained in the company of
  others are considerably harder to drop, partly because somebody notices when
  you do not appear. This does not require a formal club. A single person who
  expects you on Thursday is usually enough, and the effect holds even when
  that person has no interest in the activity itself.

  \ibgap{18}\\
  Finally, resist the temptation to measure everything. Elaborate tracking
  systems tend to collapse under their own weight, and the collapse of the
  system often takes the habit with it. A mark on a wall calendar will do. What
  you are recording is not performance but attendance, and attendance is the
  only thing that turns an intention into something you no longer have to think
  about.
\end{ibtextbody}

% ═══════════════════════ TEXT C — literary extract ═════════════════════════
\ibpagebreak
\ibtextlabel{C}
\ibtexttitle{An extract from \textit{The Long Way Round}}

\begin{ibtextbody}
  The bus did not come at six, and Mirela had known it would not. She had known
  since the previous evening, when the driver had mentioned, in the flat way he
  mentioned everything, that the road above Kesh was being resurfaced. Still she
  had walked down to the shelter at ten to, and still she had stood there in the
  cold with her bag between her feet, because standing at the shelter was what
  one did, and because to stay in the house would have meant admitting that the
  day had already gone wrong.

  At half past six she began to walk. It was eleven kilometres to the town and
  she had done it before, though not in these shoes and not with the bag. The
  road climbed for the first three and then gave up and ran flat between fields
  that were still white at the edges. She passed the Dennehy place, where a dog
  she had known since it was a puppy came to the gate and did not bark, and she
  thought, as she always thought at that gate, that she would miss the animal
  more than she would miss most of the people.

  Her mother had written to say that the town was much changed. Her mother had
  been writing to say that the town was much changed for nine years, and Mirela
  had learned to read the sentence as a kind of weather report, accurate about
  the surface and silent about everything underneath. The letters never
  mentioned the shop. They never mentioned her father's brother either, which
  amounted to the same thing.

  By eight the light had come up properly and she could see the water tower,
  which meant four kilometres, which meant she would be early. She sat on the
  wall by the crossroads and ate the bread she had brought, and a lorry went past
  and did not slow, and she found that she was not sorry. There was something in
  arriving on foot that she wanted. Let them see her come up the hill with the
  bag. Let them work out for themselves how far she had walked, and what it had
  cost her, and how little she intended to explain.
\end{ibtextbody}

\ibpagebreak
\ibdisclaimer
\ibreferences{%
  \ibrefentry{Text A}{Written for this template. Brookvale, its library and
    everyone in this article are fictitious.}
  \ibrefentry{Text B}{Written for this template.}
  \ibrefentry{Text C}{Written for this template. \textit{The Long Way Round}
    is not a real work.}}

\end{document}
